\documentclass{article}
\usepackage[utf8]{inputenc}
\usepackage{amsmath,amssymb}
\usepackage[most]{tcolorbox}
\usepackage{geometry}
\geometry{margin=2.5cm}

\newcommand{\step}[1]{\par\noindent\hangindent=3.5em\hangafter=1%
  \makebox[3.5em][l]{\textbf{#1.}}}
\newcommand{\steptag}[1]{\hfill{\small\textsf{[#1]}}}

\newtcolorbox{proofbox}[1]{
  colback=gray!5, colframe=gray!60,
  fonttitle=\bfseries, title={#1},
  breakable, enhanced,
  left=4pt, right=4pt, top=4pt, bottom=4pt,
  before skip=10pt, after skip=10pt
}

\begin{document}

\begin{proofbox}{After first fixing one global witness package W\_RF suppli...}
\step{1} After first fixing one global witness package W\_RF supplied by lem-routef-raw-factor-setting-formation, for every input (H,Phi,eta) to which that formation result applies, fix one def-routef-raw-factor-setting datum S over that same W\_RF supplied by the same result, writing the fields of (W\_RF,S) as the unqualified symbols below: Raw factor-map units: for 0 <= eta <= rho\_unit := rho\_T, max\{||tilde-Delta(I)-I||, ||tilde-Upsilon(I)-I||\} <= C\_T*eta. \steptag{validated} \\[6pt]
\step{1.1} Fix the single global header W\_RF and, for an arbitrary input (H,Phi,eta) to which lem-routef-raw-factor-setting-formation applies, fix the datum S furnished by that same external over W\_RF. Assume 0<=eta<=rho\_unit=rho\_T. By def-routef-raw-factor-setting, rho\_T is the minimum in (1.1), so rho\_T<=1/[4*(1+C\_V)]; formation gives C\_A=20+(211/8)*C\_theta>0, bar C\_E=max\{1,C\_E\}>=1, hence C\_V=bar C\_E*C\_A>=0, and therefore b:=C\_V*eta<=C\_V/[4*(1+C\_V)]<=1/4. Also C\_theta=12*(sqrt(2)-1)>=0 and C\_T=C\_theta+3*C\_V, so C\_V<=C\_T and 3*C\_V<=C\_T. \steptag{validated} \\[4pt]
\step{1.2} Delta-unit branch. Let e\_B denote the unit element of the unital C*-algebra B; by def-routef-raw-factor-setting this is the unsubscripted I forced by a B-element type, whereas I\_B denotes the identity map of B and is not used as an element. Lem-routef-raw-factor-setting-formation gives v:B->A as an extended delta-isomorphism for delta:=C\_E*epsilon\_AI(eta), with 0<=epsilon\_AI(eta)<=C\_A*eta, and def-routef-raw-factor-setting gives tilde-Delta=inclusion\_\{A subset B(H)\} composed with v. By def-extended-delta-inclusion, the n=1 amplification v is a delta-inclusion. By the registered byte-verbatim external GT-kitaev-def-delta-homomorphism, a delta-inclusion is a delta-homomorphism and its unit clause gives ||v(e\_B)-I\_A||<=delta. Since bar C\_E=max\{1,C\_E\}, delta<=bar C\_E*C\_A*eta=C\_V*eta. The unit I\_A inherited by A is I\_\{B(H)\}, and the ambient inclusion is isometric, so ||tilde-Delta(e\_B)-I\_\{B(H)\}||<=C\_V*eta<=C\_T*eta, the last inequality being node 1.1. \steptag{validated} \\[6pt]
\step{1.2.1} Dependency audit (the verifier objection is correct). The permitted text defines an extended delta-isomorphism only by bijectivity, the two-sided amplification norm bounds, and the otherwise undefined predicate delta-homomorphism; none of the registered definitions or four externals states a unit-defect inequality. This omission is mathematically material: at eta=0 take H=C, Phi=identity, tilde-Phi=identity, A=B=C, v=-identity and u=-identity. If the undefined delta-homomorphism predicate is interpreted as imposing no additional clause, then v is bijective and every amplification is an exact isometry, all explicit raw-factor norm conclusions hold, but ||v(1)-1||=2 whereas C\_V*eta=0. Thus no bridging argument from the currently permitted inputs can prove node 1.2. It becomes valid only after provisioning a permitted definition or external whose delta-homomorphism clause explicitly gives ||v(1\_B)-1\_A||<=delta; a pending definition request for precisely that clause has been filed. \steptag{validated} \\[4pt]
\step{1.2.2} Updated challenge disposition after provisioning. The earlier dependency objection no longer blocks node 1.2: the registered byte-verbatim external GT-kitaev-def-delta-homomorphism now explicitly defines a delta-homomorphism by the unit condition ||v(I)-I||<=delta and the multiplicative condition, and defines a delta-inclusion to be such a delta-homomorphism with two-sided norm bounds. Combined with def-extended-delta-inclusion at amplification n=1, formation therefore gives ||v(e\_B)-I\_A||<=delta for the correctly typed unit element e\_B of B. The former eta=0 map v=-id is not a counterexample because it violates both the unit condition and, in general, the multiplicative condition for a 0-homomorphism. The valid notation correction remains: I\_B denotes the identity map, so the amended parent uses e\_B. Consequently delta=C\_E*epsilon\_AI(eta)<=bar C\_E*C\_A*eta=C\_V*eta<=C\_T*eta proves the repaired Delta-unit branch. \steptag{validated} \\[4pt]
\step{1.2.3} The missing dependency is now explicitly provisioned as the registered byte-verbatim external GT-kitaev-def-delta-homomorphism (approximate\_algebras.tex:443-456). Formation supplies an extended delta-isomorphism v with delta=C\_E*epsilon\_AI(eta). By def-extended-delta-inclusion, its n=1 amplification v is a delta-inclusion; the external defines a delta-inclusion to be a delta-homomorphism satisfying norm bounds and defines every delta-homomorphism to obey ||v(e\_B)-I\_A||<=delta. Thus the required unit estimate follows directly, with the correctly typed unit element e\_B rather than I\_B. Since 0<=epsilon\_AI(eta)<=C\_A*eta and C\_E<=bar C\_E=max\{1,C\_E\}, delta<=bar C\_E*C\_A*eta=C\_V*eta; the inherited unit of A is I\_\{B(H)\}, the inclusion A subset B(H) preserves its norm, and validated node 1.1 gives C\_V<=C\_T. Therefore ||tilde-Delta(e\_B)-I\_\{B(H)\}||<=C\_T*eta. \steptag{validated} \\[4pt]
\step{1.3} Upsilon-unit branch. Let e\_B denote the unit element of the unital C*-algebra B; this is the unsubscripted I forced by an adjacent B-element type in def-routef-raw-factor-setting, whereas I\_B denotes the identity map of B and is not used as an element. By lem-routef-raw-factor-norms at n=1, for every X in B, ||v(X)||=||tilde-Delta(X)|| ≥ (1-C\_V*eta)||X||. Put b:=C\_V*eta. Since eta ≤ rho\_T ≤ 1/[4*(1+C\_V)], we have b ≤ C\_V/[4*(1+C\_V)] < 1/4. Because v is bijective and u=v\textasciicircum{}(-1), the lower norm bound gives ||u|| ≤ 1/(1-b). Lem-routef-raw-factor-setting-formation gives tilde-Phi\textasciicircum{}2=tilde-Phi, A=Im(tilde-Phi), A with inherited unit I, and u=v\textasciicircum{}(-1). Hence I belongs to A=Im(tilde-Phi), so idempotence gives tilde-Phi(I)=I, and def-routef-raw-factor-setting gives tilde-Upsilon(I)=u(tilde-Phi(I))=u(I). Since u(v(e\_B))=e\_B and u is linear, tilde-Upsilon(I)-e\_B=u(I)-u(v(e\_B))=u(I-v(e\_B)). The unit clause of the delta-homomorphism contained in the registered definition def-extended-delta-inclusion, applied to the extended C\_E*epsilon\_AI(eta)-isomorphism v from formation, gives ||I-v(e\_B)|| ≤ C\_E*epsilon\_AI(eta) ≤ bar C\_E*C\_A*eta=C\_V*eta. Therefore ||tilde-Upsilon(I)-e\_B|| ≤ ||u||*||I-v(e\_B)|| ≤ C\_V*eta/(1-b) ≤ 4*C\_V*eta/3 ≤ 3*C\_V*eta ≤ C\_T*eta, using b < 1/4 and C\_T=C\_theta+3*C\_V with C\_theta ≥ 0. The second unsubscripted I in the contract's expression ||tilde-Upsilon(I)-I|| has type B and is exactly e\_B. \steptag{validated} \\[6pt]
\step{1.3.1} Provisioned unit-defect bridge. The registered byte-verbatim external GT-kitaev-def-delta-homomorphism defines a delta-homomorphism v from a unital approximate Banach algebra to another by, in particular, ||v(I)-I|| <= delta. By def-extended-delta-inclusion, an extended delta-isomorphism is bijective and every amplification is a delta-inclusion, hence its n=1 amplification v is a delta-homomorphism in exactly this registered sense. Formation supplies v:B->A as an extended delta-isomorphism for delta:=C\_E*epsilon\_AI(eta). Applying the registered unit clause to the unit e\_B of B and the inherited unit I of A therefore gives ||v(e\_B)-I|| <= C\_E*epsilon\_AI(eta). Since 0<=epsilon\_AI(eta)<=C\_A*eta and C\_E<=bar C\_E:=max\{1,C\_E\}, this is at most bar C\_E*C\_A*eta=C\_V*eta. Thus the formerly unsupported estimate in node 1.3 is now an explicit consequence of a permitted byte-verbatim external plus def-extended-delta-inclusion; no conclusion is inferred from bijectivity or norm bounds alone. \steptag{validated} \\[4pt]
\end{proofbox}

\end{document}
